
\documentclass{article}
\usepackage{colortbl}
\usepackage{makecell}
\usepackage{multirow}
\usepackage{supertabular}

\begin{document}

\newcounter{utterance}

\centering \large Interaction Transcript for game `cladder', experiment `full\_v1.5\_default', episode 1290 with Qwen3.5{-}9B.
\vspace{24pt}

{ \footnotesize  \setcounter{utterance}{1}
\setlength{\tabcolsep}{0pt}
\begin{supertabular}{c@{$\;$}|p{.15\linewidth}@{}p{.15\linewidth}p{.15\linewidth}p{.15\linewidth}p{.15\linewidth}p{.15\linewidth}}
    \# & \multicolumn{2}{c}{Player} && \multicolumn{2}{c}{Game Master} \\
    \hline

    \theutterance \stepcounter{utterance}  
    & & & \multicolumn{4}{p{0.6\linewidth}}{
        \cellcolor[rgb]{0.9,0.9,0.9}{
            \makecell[{{p{\linewidth}}}]{
                \texttt{\tiny{[P1$\langle$GM]}}
                \texttt{You are an expert in causal inference. The following question is not a typical commonsense query, but rather a meticulously designed question created by a professor specializing in causal inference, intended to assess the students' mastery of the course content.} \\
\texttt{Imagine a self{-}contained, hypothetical world with only the following conditions, and without any unmentioned factors or causal relationships: Gender has a direct effect on smoking and lung cancer. Smoking has a direct effect on tar deposit. Tar deposit has a direct effect on lung cancer. Gender is unobserved. For nonsmokers, the probability of high tar deposit is 36\%. For smokers, the probability of high tar deposit is 84\%. For nonsmokers and with no tar deposit, the probability of lung cancer is 12\%. For nonsmokers and with high tar deposit, the probability of lung cancer is 44\%. For smokers and with no tar deposit, the probability of lung cancer is 15\%. For smokers and with high tar deposit, the probability of lung cancer is 50\%. The overall probability of smoking is 67\%. Will smoking increase the chance of lung cancer?} \\
\texttt{Start your answer with 'Yes' or 'No', followed by additional reasoning or evidence to support your explanation.} \\
            }
        }
    }
     \\ \\

    \theutterance \stepcounter{utterance}  
    & \multicolumn{4}{p{0.6\linewidth}}{
        \cellcolor[rgb]{0.561,0.737,0.561}{
            \makecell[{{p{\linewidth}}}]{
                \texttt{\tiny{[P1$\rangle$GM]}}
                \texttt{Yes} \\
\\ 
\texttt{**Reasoning:**} \\
\\ 
\texttt{To determine if smoking increases the chance of lung cancer, we must calculate the causal effect of smoking on lung cancer, denoted as \$P(LungCancer | Do(Smoking))\$. In this hypothetical world, the causal structure is defined as:} \\
\texttt{\$Gender \textbackslash{}rightarrow Smoking\$} \\
\texttt{\$Smoking \textbackslash{}rightarrow TarDeposit\$} \\
\texttt{\$TarDeposit \textbackslash{}rightarrow LungCancer\$} \\
\\ 
\texttt{Since \$Gender\$ is unobserved, we cannot condition on it to block the backdoor path from Smoking to Lung Cancer via Gender. However, the question asks for the total causal effect of smoking. We can calculate this by comparing the probability of lung cancer in the population where everyone is forced to smoke (\$Do(Smoking=1)\$) versus the population where everyone is forced not to smoke (\$Do(Smoking=0)\$).} \\
\\ 
\texttt{The causal effect is given by the difference: \$P(L | Do(S=1)) {-} P(L | Do(S=0))\$.} \\
\\ 
\texttt{**Step 1: Calculate \$P(L | Do(S=1))\$**} \\
\texttt{In the intervention \$Do(S=1)\$, everyone smokes. The distribution of \$TarDeposit\$ is determined solely by the smoking status (since the path from Gender to Smoking is broken by the intervention, and Smoking is the only parent of TarDeposit).} \\
\texttt{*   \$P(Tar=High | Do(S=1)) = 0.84\$} \\
\texttt{*   \$P(Tar=No | Do(S=1)) = 1 {-} 0.84 = 0.16\$} \\
\\ 
\texttt{Now, we calculate the probability of lung cancer using the provided conditional probabilities for smokers:} \\
\texttt{*   \$P(L | S=1, Tar=High) = 0.50\$} \\
\texttt{*   \$P(L | S=1, Tar=No) = 0.15\$} \\
\\ 
\texttt{\$\$P(L | Do(S=1)) = P(L | S=1, Tar=High) \textbackslash{}cdot P(Tar=High | S=1) + P(L | S=1, Tar=No) \textbackslash{}cdot P(Tar=No | S=1)\$\$} \\
\texttt{\$\$P(L | Do(S=1)) = (0.50 \textbackslash{}times 0.84) + (0.15 \textbackslash{}times 0.16)\$\$} \\
\texttt{\$\$P(L | Do(S=1)) = 0.42 + 0.024 = 0.444\$\$} \\
\\ 
\texttt{**Step 2: Calculate \$P(L | Do(S=0))\$**} \\
\texttt{In the intervention \$Do(S=0)\$, no one smokes.} \\
\texttt{*   \$P(Tar=High | Do(S=0)) = 0.36\$ (Given in the prompt)} \\
\texttt{*   \$P(Tar=No | Do(S=0)) = 1 {-}} \\
            }
        }
    }
    & & \\ \\

    \theutterance \stepcounter{utterance}  
    & & & \multicolumn{2}{p{0.3\linewidth}}{
        \cellcolor[rgb]{0.9,0.9,0.9}{
            \makecell[{{p{\linewidth}}}]{
                \texttt{\tiny{[GM$|$GM]}}
                \texttt{Yes} \\
            }
        }
    }
    & & \\ \\

    \theutterance \stepcounter{utterance}  
    & & & \multicolumn{2}{p{0.3\linewidth}}{
        \cellcolor[rgb]{0.9,0.9,0.9}{
            \makecell[{{p{\linewidth}}}]{
                \texttt{\tiny{[GM$|$GM]}}
                \texttt{yes} \\
            }
        }
    }
    & & \\ \\

    \theutterance \stepcounter{utterance}  
    & & & \multicolumn{2}{p{0.3\linewidth}}{
        \cellcolor[rgb]{0.9,0.9,0.9}{
            \makecell[{{p{\linewidth}}}]{
                \texttt{\tiny{[GM$|$GM]}}
                \texttt{game\_result = WIN} \\
            }
        }
    }
    & & \\ \\

\end{supertabular}
}

\end{document}
